\chapter{Il punteggio e la bussola}
\label{ch:score}

Il solutore del Mattone~1 risponde alla domanda ``che cosa fa il gas
in questo ugello?''. Un ottimizzatore ha bisogno di altre due cose:
un singolo numero che dica quanto è buono un ugello, e la derivata
di quel numero rispetto a ogni manopola di forma, che gli dice in
che direzione muoversi. Questo capitolo costruisce entrambi e,
soprattutto, rende entrambi verificabili --- il numero calcolandolo
due volte lungo due percorsi completamente diversi, la derivata
confrontandola con riesecuzioni a forza bruta del solutore. Aggiunge
poi un secondo strumento che non assegna un punteggio a un progetto
ma lo \emph{certifica}: un quadrante che segna zero esattamente
quando la parete dell'ugello finisce dove una parete ottimale deve
finire.

\section{Il funzionale di spinta, per due vie}

\subsection{La domanda}

Un ottimizzatore ha bisogno esattamente di due cose che il solutore
di flusso non aveva. Primo, un \emph{numero} che dica quanto è buono
un progetto --- l'\emph{obiettivo} $J$, qui la spinta in newton, la
forza assiale che l'ugello produce. Secondo, il suo
\emph{gradiente}: l'elenco delle derivate parziali di quel numero
rispetto a ciascuna manopola di progetto, cioè di quanto cambia $J$
per metro di variazione di ciascuna manopola. Il gradiente è la
bussola; punta in salita.

\subsection{Come funziona, per esteso}

La spinta è un \emph{bilancio di quantità di moto}: la seconda legge
di Newton applicata non a un corpo ma a una regione di gas. Si
traccia una \emph{superficie di controllo} --- una superficie chiusa
immaginaria che racchiude il flusso dell'ugello --- e vi si sommano
due cose: la quantità di moto che il gas trasporta attraverso la
superficie e la forza di pressione che vi spinge sopra. La somma è
la forza che la macchina sente. Nulla di fisico dipende da dove
quella superficie viene tracciata, e questa libertà viene sfruttata
due volte, calcolando $J$ su due superfici che non condividono alcun
dato:
\begin{align}
J &= \dot m\, u_e + (p_e - p_a)\, A_e
  &&\text{(via 1: leggere l'uscita),} \label{eq:route1}\\
J &= F_{\rm throat} + \int_{\rm wall} (p - p_a)\, 2\pi y \,\mathrm{d}y
  &&\text{(via 2: leggere la parete).} \label{eq:route2}
\end{align}
La via~1 chiude la superficie sul piano di uscita, dove un ugello
ideale consegna un flusso uniforme e parallelo all'asse. Lì il
bilancio è elementare: la portata $\dot m$ [kg/s] per la velocità di
uscita $u_e$ [m/s], più lo scarto fra la pressione statica di uscita
$p_e$ e la pressione ambiente $p_a$ agente sull'area di uscita
$A_e$.

La via~2 chiude la stessa superficie lungo la macchina stessa. Il
suo primo termine è il flusso di quantità di moto entrante
attraverso la \emph{gola} --- la sezione più stretta, dove il gas
raggiunge esattamente la velocità locale del suono. Una gola in
quello stato si dice \emph{bloccata}: la portata che la attraversa è
allora fissata dalla sola area di gola e dalle condizioni di camera,
e nulla a valle può aumentarla. Il carrier valuta questo flusso
sulla linea dei valori iniziali della marcia, la breve linea
attraverso la gola da cui il solutore parte. Il secondo termine è la
spinta assiale della parete: ogni anello di parete di raggio $y$ e
larghezza radiale $\mathrm{d}y$ presenta un'area assiale proiettata
$2\pi y\,\mathrm{d}y$ su cui l'eccesso di pressione $(p - p_a)$
spinge in avanti. Le due vie non condividono alcuna quantità
intermedia: la via~1 non legge mai una pressione di parete, la via~2
non legge mai lo stato di uscita.

La legge di Newton obbliga le due risposte a coincidere, quindi la
loro differenza $R_{\rm mom}$ [N] è un misuratore meccanico
d'errore: se la discretizzazione è sbagliata da qualche parte, le
due vie smettono di accordarsi. Quanto grande le è concesso essere?
Non un numero scelto a gusto. La marcia viene eseguita due volte, a
una risoluzione e a risoluzione doppia; la \emph{variazione} di
$R_{\rm mom}$ fra le due esecuzioni è la stima che lo schema fa del
proprio errore di discretizzazione. Quella variazione, moltiplicata
per il fattore di sicurezza $K_{\rm RICH} = 4$ e sommata a un
pavimento per l'arrotondamento in virgola mobile, è la banda
ammessa --- una \emph{banda di Richardson}, dal nome dell'idea di
estrapolazione a due risoluzioni da cui proviene.

Altri due strumenti completano il passo. Il \emph{test di gauge}: si
somma ovunque la stessa costante alla pressione e la spinta non deve
muoversi, perché una pressione costante integra a zero su qualunque
superficie chiusa. È un'identità delle \emph{formule}, non del
flusso, e quindi isola gli errori nella contabilità della superficie
da quelli nel flusso --- e ne ha colto uno
(Appendice~\ref{app:sliver}). E il \emph{gradiente}: poiché la
marcia è scritta in un linguaggio differenziabile, la
\emph{differenziazione automatica in modo inverso} (AD) --- la
tecnica che differenzia il programma stesso, propagando le derivate
all'indietro attraverso ogni operazione eseguita, il che è
matematicamente equivalente a risolvere le equazioni aggiunte
discrete --- restituisce $\mathrm{d}J/\mathrm{d}P$ per tutti i
parametri di progetto $P = (y_t, r_{tu}, r_{td}, \varepsilon)$ in
una sola passata all'indietro, a circa il costo di una marcia
aggiuntiva. Ciascuna delle quattro componenti è poi verificata
contro una \emph{differenza finita centrata}: si riesegue la marcia
con quella manopola spostata di $+h$ e di $-h$ e si divide la
differenza delle due spinte per $2h$. Dimezzare $h$ fornisce la
seconda stima con cui si costruisce la banda di Richardson.

\begin{figure}[htbp]
\centering
\includegraphics[width=0.82\textwidth]{sch_two_routes.png}
\caption{Due vie indipendenti verso un'unica spinta, disegnate nelle
coordinate proprie dell'ugello ($x$ [m] lungo l'asse, $y$ [m]
raggio). La via~1 (a destra) chiude la superficie di controllo sul
piano di uscita e richiede solo lo stato uniforme di uscita: portata
per velocità di uscita più $(p_e-p_a)A_e$. La via~2 (a sinistra e
lungo la parete) la chiude sulla macchina: flusso di quantità di
moto attraverso la gola bloccata più la spinta assiale di pressione
$(p-p_a)\,2\pi y\,\mathrm{d}y$ di ogni anello di parete. Si guardi a
ciò che le due vie \emph{non} condividono: nessuna quantità è usata
da entrambe, quindi la loro differenza $R_{\rm mom}$ [N] è un
autentico misuratore d'errore e non un'identità. Il cerchio rosso
segna la scheggia di superficie lasciata inizialmente aperta per
sbaglio, la cui storia è raccontata
nell'Appendice~\ref{app:sliver}.}
\label{fig:tworoutes}
\end{figure}

\subsection{Che cosa è stato costruito}

\carrier{validation/a1\_thrust\_functional.py}: entrambe le vie, il
residuo di quantità di moto $R_{\rm mom}$ con la sua banda di
Richardson, il test di gauge, $\mathrm{d}J/\mathrm{d}P$ in modo
inverso verificato manopola per manopola contro differenze finite
centrate, e tre rigettatori --- versioni deliberatamente corrotte
del calcolo (termine pressione--area eliminato, superficie non
chiusa, flusso corrotto) che la suite deve cogliere.

\subsection{Risultati di riferimento}

\textbf{8/8 PASS.} $J = 1.2552079\times10^{8}$~N; $R_{\rm mom} =
+2.85\times10^{3}$ contro una banda derivata di
$4.01\times10^{4}$ ($3.2\times10^{-4}$ di $J$), quindi le due vie
concordano a circa tre parti su diecimila della spinta; il rapporto
grossolana/fine dei residui pari a 4.52 è il $\approx 4$ atteso da
uno schema del secondo ordine quando la risoluzione viene
raddoppiata, il che conferma che ciò che resta è errore di
discretizzazione e non un errore concettuale. Scarto di gauge
$3.35\times10^{-8}$ contro $1.78\times10^{-6}$ ammesso: aggiungere
una pressione costante davvero non cambia nulla. Tutte e quattro le
componenti del gradiente cadono dentro le rispettive bande da
differenze finite ($\mathrm{d}J/\mathrm{d}y_t$ concorda a 9 cifre
significative), e tutti i rigettatori scattano.

Anche i segni del gradiente si comportano come la teoria prevede:
$\mathrm{d}J/\mathrm{d}\varepsilon > 0$ contro il vuoto, cioè
quando nulla spinge in senso contrario conviene sempre costruire un
ugello più lungo; e $\mathrm{d}J/\mathrm{d}r_{td}$ è minuscolo,
cioè la curvatura del breve raccordo appena a valle della gola conta
pochissimo. Un difetto reale è stato colto e corretto lungo la
strada: Appendice~\ref{app:sliver}.

\section{Il misuratore di ottimalità}
\label{sec:corner}

\subsection{La domanda}

Un punteggio e una bussola trovano progetti migliori. Che cosa
\emph{certifica} che un dato progetto soddisfa già la condizione che
un ottimo deve soddisfare?

\subsection{Come funziona, per esteso}

Scegliere la migliore parete di ugello non è un problema di trovare
il valore migliore di pochi numeri, ma di trovare la migliore
\emph{curva}; il ramo della matematica che se ne occupa è il
\emph{calcolo delle variazioni}, che si chiede che cosa debba essere
vero di una curva se nessuna sua piccola deformazione può migliorare
l'obiettivo. Rao~\cite{rao} risolse quel problema per un ugello e
trovò, fra le condizioni, un'affermazione scalare esatta che deve
valere al labbro --- il punto in cui la parete si ferma. È una
condizione \emph{terminale} (o di trasversalità):
\begin{equation}
R_{\rm corner}(p_a) \;=\; (p - p_a)\;-\;\tfrac{1}{2}\rho W^{2}
\sin(2\theta)\tan\alpha .
\label{eq:corner}
\end{equation}
La si legga come il bilancio dei due effetti dello spostare il
labbro di un passo più a valle. Il primo termine è il guadagno
diretto: un anello di parete in più, che conviene costruire finché
la pressione su di esso supera l'ambiente. Il secondo è il prezzo:
la parete aggiuntiva costringe il flusso a continuare a ruotare, e
il costo di quella rotazione è scritto attraverso l'angolo locale
del flusso $\theta$ (la direzione della velocità rispetto all'asse)
e l'angolo di Mach $\alpha$ (l'inclinazione delle linee lungo cui
viaggiano le perturbazioni supersoniche,
Capitolo~\ref{ch:moc}). $R_{\rm corner} = 0$ significa ``la parete
si ferma esattamente nel punto giusto'': un anello in più
guadagnerebbe quanto costa. Il misuratore valuta questa quantità dai
dati di labbro della marcia stessa, così che qualunque progetto il
motore produca possa essere certificato e non soltanto valutato.
Mediata sul ciclo del motore diventa la condizione di labbro che
l'ugello di un motore a detonazione rotante deve soddisfare
(affermazione T7 di~\cite{theory}), che il Capitolo~\ref{ch:cycle}
mette su un misuratore dedicato.

Lo strumento compagno è il \emph{primo integrale}
\begin{equation}
f_2 \;=\; \frac{W\cos(\theta-\alpha)}{\cos\alpha} \;=\; -\lambda_2 ,
\label{eq:f2}
\end{equation}
dove ``primo integrale'' significa una combinazione di grandezze di
flusso che le equazioni del moto obbligano a restare costante lungo
una particolare curva, benché ciascun ingrediente vari. La teoria di
Rao rende $f_2$ costante lungo la \emph{caratteristica terminale} di
un ugello ottimale --- l'ultima delle due famiglie di linee lungo
cui viaggia l'informazione supersonica, quella che chiude il campo
di moto al labbro. Il suo valore costante è $-\lambda_2$, un
\emph{moltiplicatore di Lagrange}: il numero che dà il prezzo di un
vincolo in un'ottimizzazione vincolata, qui indicando quanta spinta
si guadagnerebbe per ogni unità di portata in più concessa. Lungo
qualunque altra curva --- in particolare lungo la parete, che è una
\emph{linea di corrente}, cioè una linea ovunque tangente alla
velocità e dunque il percorso seguito da una particella di gas ---
$f_2$ \emph{non} è costante. Questo fornisce il controllo negativo
naturale: un'implementazione che facesse apparire $f_2$ costante
ovunque sarebbe rotta, e va colta.

\subsection{Che cosa è stato costruito}

\carrier{validation/a1\_corner\_instrument.py}: il misuratore di
spigolo e lo strumento $f_2$, verificati contro tre risposte note
sul caso dell'ugello ideale (uscita assiale, ambiente adattato,
vuoto), con il rigettatore di non costanza sulla parete.

\subsection{Risultati di riferimento}

\textbf{5/5 PASS} (82~s). L'angolo del flusso al labbro è zero di
macchina ($\theta = -10^{-16}$~rad: l'uscita è davvero assiale, come
l'ugello ideale assume). Con l'ambiente posto uguale alla pressione
di uscita il misuratore segna $3.7\times10^{-10}$ contro 42.8~Pa
ammessi --- l'ugello è certificato stazionario al suo ambiente di
progetto. Contro il vuoto segna esattamente $+p_e$, cioè
``allungami'', l'affermazione che nessun ugello finito è ottimale
nel vuoto, ora su un quadrante; e il suo segno concorda con un
$\mathrm{d}J/\mathrm{d}\varepsilon$ ricalcolato in modo
indipendente, così misuratore e gradiente raccontano la stessa
storia. Il primo integrale $f_2$ al labbro è uguale alla velocità di
uscita a meno di $4.6\times10^{-13}$~m/s, mentre lungo la parete
varia di 1173~m/s --- il rigettatore di non costanza, che scatta
come previsto.
